%=====================================================================
\chapter{The Boundary with the Neighbouring Courses}\label{app:boundary}
%=====================================================================
\normalfigures

Artificial Intelligence is the gateway of this cluster: it is the published
prerequisite for \emph{Machine Learning}, \emph{Programming for Artificial
Intelligence}, \emph{Knowledge Representation and Reasoning} and
\emph{Introduction to Data Science}. That makes its boundaries unusually
consequential --- a topic this course teaches badly is a topic four other courses
inherit badly, and a topic it teaches twice is a week nobody gets back.

This appendix states each boundary explicitly: what this course teaches, what it
deliberately does not, and which course picks it up.

\section{Why This Course Is Symbolic}

The decision that shapes everything else: \emph{this course teaches the agent that
is told the rules}. It is given the tasks and their dependencies, the constraints,
the operators, the facts and the probabilities, and its job is to work out what to
do or what follows. It is not given examples and asked to induce a rule --- that is
the next course.

Three reasons, none of them nostalgia.

\begin{enumerate}[leftmargin=2.2em, itemsep=3pt, topsep=4pt]
  \item \textbf{This is where guarantees live.} A constraint solver can prove that
        no engineer is double-booked (\chref{csp}). No learned model can prove
        anything of the kind, and for a class of requirements --- safety,
        regulation, contractual obligation --- a proof is what is actually
        required.
  \item \textbf{This is where explanations live.} A rule engine prints the chain of
        facts and inferences that produced its conclusion (\chref{expert}). That
        chain is the artefact an auditor accepts, and \chref{responsible} argues
        that it is an accountability instrument rather than a convenience.
  \item \textbf{This is what you have before you have data.} Every project begins
        with no examples and a domain expert who knows the rules. Symbolic methods
        are the ones that work on day one.
\end{enumerate}

Semester~3 is also simply too early for the statistics that machine learning
needs, and the HEC 2023 prerequisite for this course --- Object Oriented
Programming --- confirms it.

\section{Machine Learning (Semester 6)}

\emph{Machine Learning} has \emph{Artificial Intelligence} as its published HEC
2023 prerequisite. It owns learning entirely.

\rowcolors{2}{white}{lightbg}
\begin{longtable}{@{}L{4.6cm}L{4.6cm}L{5.4cm}@{}}
\toprule
\rowcolor{primary!15}
\textbf{Topic} & \textbf{Here} & \textbf{There} \\
\midrule
\endhead
Learning paradigms & \chref{learning}: named, one delivery example each, one
figure & Chapters 1--4: the formal setting, hypothesis spaces, inductive bias \\
Neural networks & \chref{learning}: the perceptron, the learning rule by hand,
and the exclusive-or limit. The course stops there & Chapter 13: multi-layer
networks, the delta rule, backpropagation, activation functions \\
Markov decision processes, reinforcement learning & Not taught. The
agent--reward loop appears once as a figure in \chref{learning} & Chapters 18--20:
hidden Markov models, MDPs, Monte Carlo, temporal difference, Q-learning \\
Naive Bayes & Not taught. \chref{bayesnets} does \emph{inference} in a network
whose probabilities are given & Chapter 10: \emph{learning} those probabilities
from data, with smoothing \\
Evaluation and generalisation & Not taught at all & Chapters 7--8: the hinge of
that course \\
\bottomrule
\end{longtable}

\begin{conceptbox}
The clean way to state this boundary in one sentence: \emph{this course is given
the probabilities and computes with them; Machine Learning estimates the
probabilities from data.} \chref{bayesnets} and Machine Learning's Chapter~10 look
superficially similar and are doing opposite halves of the job.
\end{conceptbox}

\section{Knowledge Representation and Reasoning (elective)}

The boundary most easily overlooked, because the overlap is largest.
\emph{Knowledge Representation and Reasoning} is an HEC 2023 elective, 3 (3-0),
with \emph{Artificial Intelligence} as its prerequisite. Its published outline
includes propositional logic, first-order logic, Horn clauses, description logic,
reasoning using description logic, forward and backward chaining, semantic
networks, ontologies and ontology languages, logical agents, planning, rule-based
knowledge representation, reasoning under uncertainty, Bayesian networks, fuzzy
logic, Markov models, commonsense reasoning and explainable AI.

Almost all of Parts~IV and~V of this handout appear in that list. That is not an
error in either document: this course teaches those topics \emph{once, at first
pass}, because inference cannot be taught without a representation to reason over,
and because the HEC 2017 outline for \emph{Artificial Intelligence} names
propositional logic, first-order logic, uncertainty and fuzzy logic directly.

\rowcolors{2}{white}{lightbg}
\begin{longtable}{@{}L{4.6cm}L{4.6cm}L{5.4cm}@{}}
\toprule
\rowcolor{primary!15}
\textbf{Topic} & \textbf{Here (first pass)} & \textbf{There (in depth)} \\
\midrule
\endhead
Propositional and first-order logic & \chref{propositional},
\chref{fol}: semantics, CNF, DPLL, unification, resolution refutation &
Completeness proofs, Horn clauses as a fragment, logical agents \\
Description logic & Not taught & The core of that course: concept languages,
subsumption, reasoning over terminologies \\
Ontologies & \chref{kr}: what an ontology is and what it buys, in one section &
Ontology languages, engineering methodologies, reasoning over them \\
Chaining and rule bases & \chref{expert}: a working engine, forward and backward,
with explanation & Rule-based knowledge representation as a formalism \\
Bayesian networks & \chref{bayesnets}: inference in a given network & Network
semantics as a knowledge representation \\
Commonsense reasoning, explainable AI & Not taught. \chref{responsible} treats
explanation as an accountability artefact, not as a formalism & Both, as
formalisms \\
\bottomrule
\end{longtable}

\section{Data Science (Semester 7) and the 2025 Arrivals}

\emph{Introduction to Data Science} also lists \emph{Artificial Intelligence} as
its HEC 2023 prerequisite, and the departmental \emph{Data Science Fall 2026}
handout sits at the end of the sequence.

\rowcolors{2}{white}{lightbg}
\begin{longtable}{@{}L{4.6cm}L{4.6cm}L{5.4cm}@{}}
\toprule
\rowcolor{primary!15}
\textbf{Topic} & \textbf{Here} & \textbf{There} \\
\midrule
\endhead
Natural language processing & \chref{nlp}: the \emph{symbolic} half ---
grammars, chart parsing, pattern substitution, n-grams & Data Science Chapter 20:
the neural half --- embeddings, transformers \\
Generative and agentic AI & \chref{trends}: named as trends, with one offline
tool-using agent built from \chref{expert}'s rule engine & Data Science
Chapters 21--22 \\
Deep learning & Not taught & Data Science Chapters 17--18 \\
Ethics & \chref{responsible}: \emph{can this system's reasoning be trusted and
audited?} & Data Science Chapter 28: privacy, security and duty to the data
subject. Machine Learning Chapter 21: \emph{is the learned model fair?} \\
\bottomrule
\end{longtable}

Ethics is the one place where three courses legitimately own the same word. Each
asks a different question, and the three questions are not interchangeable; the
progression one-pager distributed with this handout prints all three side by side.

\section{Programming for Artificial Intelligence}

This HEC 2023 course also has \emph{Artificial Intelligence} as its prerequisite
and covers AI programming languages, production-quality code, and the numerical
Python stack (NumPy, pandas, Matplotlib). \chref{python} here is deliberately
narrow by comparison: it fixes the four data structures this course's algorithms
need and nothing else. It is not a Python course and does not attempt to be one.

\section{The Summer 2026 Handout}

\texttt{Summer\_2026/AI\_Handout.pdf} is a 24-page overview written for an
eight-week session. It covers the same ground at a fraction of the depth and omits,
among other things, arc consistency, resolution refutation, variable elimination,
Mamdani inference and the entire analysis framework of \chref{classic}. It is a
revision summary and a source of examination practice; it is not a substitute for
this volume, and a student who has read only it has not met the HEC outline.

\begin{alertbox}[If you are planning the degree]
Two failure modes to watch for, both observed in practice.

\smallskip
\textbf{Teaching machine learning here ``because students expect it''.} It
displaces the symbolic material, which no other core course covers, and it
duplicates a course three semesters later that will teach it properly with the
statistics to support it. The students who enjoyed the taste arrive at
\emph{Machine Learning} believing they have already done it.

\smallskip
\textbf{Assuming the elective will cover the logic.} \emph{Knowledge
Representation and Reasoning} is an \emph{elective}. Most students will never take
it. If this course does not teach propositional and first-order inference, most
graduates of the programme will never meet them --- and both HEC outlines for
\emph{Artificial Intelligence} require them.
\end{alertbox}
